\section{Strategies for Subsets Mixing and Upsampling with \dolma}
\label{sec:mixtures}


Like the pretraining corpora of nearly every large-scale language model, \dolma is a multi-source dataset. 
Training on \dolma thus requires a mixing strategy that determines how much data from each source to include, and potentially which sources to upsample. 
Like other multi-source corpora (e.g., ROOTS~\citep{Laurenccon2023TheBR}, the Pile~\citep{Gao2020ThePA}, RedPajama v1~\citep{together2023redpajama}),\footnote{RedPajama v1 was a reproduction of the multi-source corpus used in LLaMA~\citep{Touvron2023LLaMAOA}. RedPajama v2~\citep{redpajama2} focuses solely on Common Crawl and is thus single-source.} \dolma does not prescribe a single mixing strategy. 
We refer the reader to \citet{Rae2021ScalingLM} for an example of how one might programmatically search over mixing configurations to maximize performance. 
Here, we perform mixing experiments as an opportunity to answer some research questions about how different data sources interact. 
We use the same ablation setup described in \autoref{sec:creation}.



\paragraph{How much code is important for pretraining?} It is common practice for language models to be pretrained on some amount of code, even if code generation is not the intended task.
Some research has suggested that mixing code into training over plain text documents improves performance on reasoning tasks~\citep{madaan-etal-2022-language}.
We investigate whether this observation holds for models trained on \dolma, and if so, how much code is needed?

\begin{table*}[h]
    \centering
   
    \small
    \begin{tabular}{llll}
        \toprule
        \textbf{Dataset} & \textbf{0\% Code} & \textbf{5\% Code} & \textbf{15\% Code} \\
        \midrule
	bAbI (ICL) & 0.0~~~\textpm~0.0 & 8.8~~~\textpm~0.9 & 10.1 \textpm~2.8 \\
	WebNLG (ICL) & 16.8 \textpm~1.1 & 19.3 \textpm~1.1 & 22.0 \textpm~1.3 \\
	GSM8K (FT) & 0.0~~~\textpm~0.0 & 0.0~~~\textpm~0.0 & 0.0~~~\textpm~0.0 \\
	GSM8K+PAL (FT) & 11.8 \textpm~0.8 & 14.2 \textpm~1.3 & 14.7 \textpm~0.9 \\
    \bottomrule
    
    \end{tabular}
    \caption{
    Performance of three models pre-trained with increasing amounts of code on three datasets, across 5 random seeds. 
    We measure exact match for bAbI and GSM8K, and Rouge-2 for WebNLG.
    }
    \label{tab:code-reasoning}
\end{table*}




We create three mixtures from the C4 and Stack subsets containing 0\%, 5\% and 15\% of code data.
On each, we train a 1B model. 
We evaluate these models on three different reasoning tasks: bAbI~\citep{Weston2015TowardsAQ}, WebNLG \cite{gardent-etal-2017-creating} and GSM8k \cite{Cobbe2021TrainingVT}. For the first two tasks, we follow the experimental setup of \citet{muennighoff2023scaling} and evaluate each model in an ICL setup with a changing number of demonstrations (0-5) across 5 random seeds. \citet{muennighoff2023scaling} show that adding code to pre-training data improves ICL performance on bAbI and WebNLG and they suggest that code improves long-range state-tracking capabilities. Our experiments, as shown in Table~\ref{tab:code-reasoning}, corroborate these findings: while the C4-only model fails on all bAbI tasks, adding code improves performance, with a similar trend for WebNLG. 

On the more difficult GSM8k benchmark, all models failed to get any correct answer in an ICL setup, and even when fine-tuning the models on the entire training set. However, we find that by fine-tuning on program-aided output, where questions are solved by writing Python snippets as described in \cite{gao2022pal}, code models outperform the C4-only model. These results show that models pre-trained on code can leverage code generation to answer challenging reasoning tasks even when the original task does not directly involve code.


\paragraph{Evaluating mixing strategies for pretraining on \dolma}
While \dolma does not prescribe a specific source mixture, we analyze some commonly used strategies\footnote{
    We did not include any social data in these mixes as it was not ready at the time of this experiment.
} and compare their effect using the Paloma evaluation suite~\citep{paloma}.
Specifically, we present and evaluate four possible data mixtures in \autoref{table:mix}. 




\begin{table*}[h]
\centering
\small
\renewcommand{\arraystretch}{1.2}
\begin{tabular}{lm{19em}cc}
    \toprule
    \textbf{Mix Name} &
    \textbf{Description} &
    \textbf{Sampling} &
    \textbf{Proportion} \\
    \midrule
    \textbf{
        Na\"ive
    } & {
        Sample each source in \autoref{tab:statistics} equally.
    }  &  {
        \renewcommand{\arraystretch}{1}
        \scriptsize
        \begin{tabular}{@{}lr@{}}
        {\tiny\dolmaWeb}~Web & 100\% \\
        {\tiny\dolmaCode}~Code & 100\% \\
        {\tiny\dolmaRefs~\dolmaPapers}~Ref. & 100\% \\
        {\tiny\dolmaBooks}~Books & 100\% \\
        \end{tabular}
    }  & {
        \renewcommand{\arraystretch}{1}
        \scriptsize
        \begin{tabular}{@{}lr@{}}
        {\tiny\dolmaWeb}~Web & 83.5\% \\
        {\tiny\dolmaCode}~Code & 13.8\% \\
        {\tiny\dolmaRefs~\dolmaPapers}~Ref. & 2.5\% \\
        {\tiny\dolmaBooks}~Books & 0.2\% \\
        \end{tabular}
        \vspace{1em}
    } \\
    \textbf{Web Only} &
    {
        Similar to~\citet{ayoola-etal-2022-refined}, we test a mixture that only uses web data.
        
        
    }  &  {
        \renewcommand{\arraystretch}{1}
        \scriptsize
        \begin{tabular}{@{}lr@{}}
        {\tiny\dolmaWeb}~Web & 100\% \\
        {\tiny\dolmaCode}~Code & 0\% \\
        {\tiny\dolmaRefs~\dolmaPapers}~Ref. & 0\% \\
        {\tiny\dolmaBooks}~Books & 0\% \\
        \end{tabular}
    }  & {
        \renewcommand{\arraystretch}{1}
        \scriptsize
        \begin{tabular}{@{}lr@{}}
        {\tiny\dolmaWeb}~Web & 100\% \\
        {\tiny\dolmaCode}~Code & 0\% \\
        {\tiny\dolmaRefs~\dolmaPapers}~Ref. & 0\% \\
        {\tiny\dolmaBooks}~Books & 0\% \\
        \end{tabular}
    }  \\
    \textbf{
        Reference+
    } & {
        It is common practice to upsamole knowledge-intensive documents when composing training mixture. 
        In our case, we upsample the PeS2o papers, Wikipedia, Wikibooks, and Gutenberg books subsets by 2x.
    }&  {
        \renewcommand{\arraystretch}{1}
        \scriptsize
        \begin{tabular}{@{}lr@{}}
        {\tiny\dolmaWeb}~Web & 100\% \\
        {\tiny\dolmaCode}~Code & 100\% \\
        {\tiny\dolmaRefs~\dolmaPapers}~Ref. & 200\% \\
        {\tiny\dolmaBooks}~Books & 200\% \\
        \end{tabular}
    }  & {
        \renewcommand{\arraystretch}{1}
        \scriptsize
        \begin{tabular}{@{}lr@{}}
        {\tiny\dolmaWeb}~Web & 81.2\% \\
        {\tiny\dolmaCode}~Code & 13.5\% \\
        {\tiny\dolmaRefs~\dolmaPapers}~Ref. & 4.9\% \\
        {\tiny\dolmaBooks}~Books & 0.4\% \\
        \end{tabular}
        \vspace{1em}
    }  \\
    \textbf{
        Gopher-like
    } & {
        Following~\citet{Rae2021ScalingLM}, we create a mix that is heavily biased towards reference material. 
        As we do not have access to the same sources, an exact replication of their mix is not possible.
    }&  {
        \renewcommand{\arraystretch}{1}
        \scriptsize
        \begin{tabular}{@{}lr@{}}
        {\tiny\dolmaWeb}~Web & 17\% \\
        {\tiny\dolmaCode}~Code & 8\% \\
        {\tiny\dolmaRefs~\dolmaPapers}~Ref. & 200\% \\
        {\tiny\dolmaBooks}~Books & 200\% \\
        \end{tabular}
    }  & {
        \renewcommand{\arraystretch}{1}
        \scriptsize
        \begin{tabular}{@{}lr@{}}
        {\tiny\dolmaWeb}~Web & 68.4\% \\
        {\tiny\dolmaCode}~Code & 5.4\% \\
        {\tiny\dolmaRefs~\dolmaPapers}~Ref. & 24.2\% \\
        {\tiny\dolmaBooks}~Books & 2.0\% \\
        \end{tabular}
    }  \\
    \bottomrule
\end{tabular}
\vspace{1em}
\caption{
    Overview of the mixtures and their composition. 
}
\label{table:mix}
\end{table*}


\begin{figure*}[h]
    \centering
    \begin{subfigure}{.31\textwidth}
    \includegraphics[width=\linewidth]{experiments/ablations_dolma_mix/ppl/c4_100_domains.pdf}
    
    \label{fig:dolma_mix:c4_100}
    \end{subfigure}
    \quad
    \begin{subfigure}{.31\textwidth}
    \includegraphics[width=\linewidth]{experiments/ablations_dolma_mix/code/humaneval.pdf}
    \end{subfigure}
    \quad
    \begin{subfigure}{.31\textwidth}
    \includegraphics[width=\linewidth]{experiments/ablations_dolma_mix/ppl/m2d2_s2orc.pdf}
    \end{subfigure}
    \caption{
    1B model ablations for different proportions of \dolma data. All mixture perform similarly on web data (left), while excluding code increases perplexity on code datasets (center). 
    Finally, increasing reference material by upsampling papers and Wikipedia yields lower perplexity on S2ORC (right). 
    Overall, source distribution is linked to downstream capabilities; thus, \dolma users should sample subsets according to their needs.
}%
    \label{fig:dolma_mix}
\end{figure*}

We show results of mixtures in \autoref{fig:dolma_mix}. 
Overall, we observe that the different mixtures have an effect on the ability of resulting models to capture specific subdomains. 
All mixtures show similar perplexity scores on pages sampled from 100 domains from C4 (\autoref{fig:dolma_mix}, left), indicating their general effectiveness at modeling web documents.  
On the other hand, we note how models struggle to model specialized domains unless they are exposed to them. 
As an example, a model trained on the \textit{Web-only} mix struggles to represent data in the code domain (\autoref{fig:dolma_mix}, center, HumanEval).
Finally, we use results on the S2ORC subset of M2D2, which consists of academic papers, to illustrate how different data mixtures affect perplexity. 
As is it the case with code, \textit{Web-only} model exhibits higer perplexity due to domain mismatch. 
On the other hand, models trained on \textit{Reference+} and \textit{Gopher-like} mixes achieve lower perplexity than the model trained on the \textit{Na\"ive} mix, due to more in-domain content. 
However, we note that, despite significant differences in the amount of academic papers between \textit{Reference+} and \textit{Gopher-like} (4.9\% vs 24.2\%), they achieve nearly identical results, suggesting that even a relatively small percentage of in-domain data is sufficient to achieve good domain fit. 
